\documentclass[12pt]{article}
\usepackage[utf8]{inputenc}
\usepackage[T1]{fontenc}
\usepackage[margin=1in]{geometry}
\usepackage{setspace}
\usepackage{xcolor}
\usepackage[colorlinks=true, urlcolor=blue!60!black, linkcolor=blue!60!black]{hyperref}
\setlength{\parindent}{0em}
\setlength{\parskip}{0.6em}
\renewcommand{\baselinestretch}{1.0}
\begin{document}
\begin{center}
{\LARGE \textit{Haaland v. Brackeen}, 599 U.S. 255 (2023)\par}
\vspace{0.5em}
{\large Case Brief --- Two-Page Edition\par}
\end{center}

\section{Memory Jogger}
ICWA survives: it is within Congress's Indian-affairs power, does not commandeer the States, and its equal protection challengers lacked standing (\href{https://supreme.justia.com/cases/federal/us/599/21-376/}{599 U.S. 255}).

\section{Detailed Case Facts}
ICWA requires ``any party'' seeking to remove an Indian child to show ``active efforts'' to keep the family together and sets adoptive preferences for extended family, tribal members, and ``other Indian families.'' Non-Indian foster and adoptive families, including the Brackeens, challenged it.

\section{Procedural History}
The families and Texas sued federal officials in the Northern District of Texas and won; the en banc Fifth Circuit upheld Congress's power but struck several provisions under the Tenth Amendment.

\section{Judicial Votes}
7--2. Barrett wrote for the Court; Gorsuch and Kavanaugh concurred; Thomas and Alito dissented.

\section{Holding}
ICWA is valid Article~I legislation with no ``family law carveout''; its rules apply evenhandedly or preempt state law in state courts; and the challengers lacked standing on equal protection.

\section{Analysis of Opinions}
\textbf{Majority:} Congress's power rests on the Indian Commerce Clause, treaties, structure, and the trust relationship. Section 1912 binds ``any party''; \S~1915 imposes no search duty on state agencies; recordkeeping is ancillary. The individuals' injury was not redressable, and Texas has no equal protection rights. \textbf{Gorsuch:} an originalist ``Indian-law bargain'' protects tribal sovereignty. \textbf{Kavanaugh:} equal protection ``is serious.'' \textbf{Thomas and Alito:} plenary power is unmoored from the text.

\section{Examples: Future Applications}
\textbf{Same-Side:} an adoption agency suing the Secretary lacks redressability. \textbf{Opposite-Side:} a search duty imposed only on state agencies likely commandeers. \textbf{Fence-Sitter:} an equal protection challenge raised in a state adoption case.

\section{Critique}
The Court declined to define the limits of a precedent it called ``unwieldy,'' and critics say the ``any party'' test elevates form over substance. Gorsuch and Thomas reached opposite originalist answers.

\section{Key Quotations}
\begin{quote}
``It is a federal court's judgment, not its opinion, that remedies an injury; thus it is the judgment, not the opinion, that demonstrates redressability. The individual petitioners can hope for nothing more than an opinion, so they cannot satisfy Article III.''
\end{quote}
\noindent --- Justice Barrett, Opinion of the Court, 599 U.S. 255.

\end{document}
